% Oversells «Oshpazlik kitobi» shabloni (XeLaTeX yoki LuaLaTeX).
% \taom{nomi}{vaqt}{porsiya}{rasm} — taom boshi (rasm bilan);
% \begin{retsept} masalliqlar \tayyorlash qadamlar \end{retsept} — ikki ustun.
% Har bir masalliq: \m{miqdor}{nima}; qadamlar: \item. Koʻrinish sahifa.sty da.
\documentclass[11pt,twoside]{article}
\usepackage[paperwidth=190mm,paperheight=240mm]{geometry}
\usepackage{fontspec}
\setmainfont{Noto Serif}
\usepackage{babel}
\babelprovide[import=uz-Latn, main]{uzbek}
\def\uzbekhyphenmins{22}
\usepackage{enumitem}
\raggedbottom
\usepackage{sahifa} % Oversells sahifa dizayni (sahifa.sty)

\newcommand\taom[4]{\clearpage\section*{#1}\markboth{#1}{}%
  \begin{uslub}{malumot}Tayyorlash vaqti: #2 \quad·\quad #3 kishilik\end{uslub}%
  \rasmramka[en=0.8\linewidth]{#4}{}}
% Ingredients on the left, the method on the right.
\makeatletter
\NewDocumentEnvironment{retsept}{+b}{\retsept@split#1\retsept@end}{}
\def\retsept@split#1\tayyorlash#2\retsept@end{%
  \par\noindent\begin{minipage}[t]{0.36\linewidth}
    \begin{uslub}{bolim}Masalliqlar\end{uslub}
    \begin{itemize}[leftmargin=*,itemsep=1pt,label=–]#1\end{itemize}
  \end{minipage}\hfill
  \begin{minipage}[t]{0.58\linewidth}
    \begin{uslub}{bolim}Tayyorlash\end{uslub}
    \begin{enumerate}[leftmargin=*,itemsep=3pt]#2\end{enumerate}
  \end{minipage}\par}
\makeatother
\newcommand\m[2]{\item \matnuslubi{miqdor}{#1} #2}

\begin{document}

\taom{Palov}{1 soat 30 daqiqa}{6}{rasmlar/osh.png}
\begin{retsept}
\m{1 kg}{guruch (devzira)}
\m{1 kg}{qoʻy goʻshti}
\m{1 kg}{sabzi}
\m{2 ta}{piyoz}
\m{200 ml}{oʻsimlik yogʻi}
\m{1 bosh}{sarimsoq}
\m{1 choy qoshiq}{zira, tuz}
\tayyorlash
\item Guruchni bir necha marta yuvib, iliq suvda ivitib qoʻying.
\item Qozonda yogʻni qizdirib, piyozni qizarguncha qovuring.
\item Goʻshtni qoʻshib, har tomonini qizartiring.
\item Somoncha toʻgʻralgan sabzini solib, 10 daqiqa qovuring; zira va tuz qoʻshing.
\item Suv quyib, 30 daqiqa past olovda qaynating (zirvak).
\item Guruchni tekis yoyib, suvini guruchdan 1,5 sm yuqori qilib quying.
\item Suv singgach, sarimsoqni oʻrtaga botiring, guruchni toʻplab, qopqoq ostida 25 daqiqa dimlang.
\end{retsept}

\begin{uslub}{maslahat}
Maslahat: guruchni dimlashdan oldin choʻmich dastasi bilan bir necha joyidan teshib qoʻying — bugʻ tekis tarqaladi.
\end{uslub}

\taom{Achchiq-chuchuk salat}{15 daqiqa}{4}{rasmlar/salat.png}
\begin{retsept}
\m{4 ta}{pomidor}
\m{1 ta}{piyoz}
\m{1 ta}{achchiq qalampir}
\m{1 bogʻ}{rayhon}
\m{}{tuz, qora murch}
\tayyorlash
\item Pomidorni yupqa yarim halqa qilib toʻgʻrang.
\item Piyozni juda yupqa toʻgʻrab, sovuq suvda 5 daqiqa ushlab turing.
\item Hammasini aralashtirib, tuz, murch va maydalangan rayhon seping.
\end{retsept}

\end{document}
